% Review
\subsection{Codecademy}

Druhou analyzovanou aplikací je aplikace Codecademy. Codecademy je jednou z nejpopulárnějších platforem na výuku programování a softwarového inženýrství. Nabízí přes 800 kurzů a má napříč platformami přes 40 milionů uživatelů. Codecademy je k dispozici jako webová aplikace a dále také jako aplikace Codecademy Go, která je k dispozici na platformách Android a iOS. 

Mobilní aplikace Codecademy Go obsahuje pouze jednoduché funkce pro opakování a procvičování látky. Hlavní náplň kurzů se pak nachází ve webové aplikaci Codecademy. Proto se v této analýze primárně zaměřím na webovou aplikaci, nicméně zahrnu do ní i poznatky ze zmíněné mobilní aplikace.  

% Review
\subsubsection{Vzhled a uživatelská zkušenost}

% Odstavec: Obecný dojem (komplexita/přehlednost UI, paleta barev, typografie)
Codecademy na první pohled působí mnohem komplexnějším dojmem než dříve analyzovaná aplikace Mimo. Design uživatelského rozhraní je spíše minimalistický, nicméně místy může stále působit příliš složitě. V aplikaci je jako barva pozadí využívána primárně světle růžová barva. Při učení je pak pro pozadí využívána bílá, černá nebo tmavě šedá barva. Pro tlačítka a další důležité prvky rozhraní pak aplikace využívá vysoce kontrastní barvy jako například modrou nebo žlutou. Co se typografie týče, Codecademy využívá napříč celou aplikací, s výjimkou bloků kódu, bezpatkové písmo Apercu.

% Odstavec: Rozložení prvků na obrazovce (lekce v kurzu, učení)
Rozložení prvků na obrazovce je poměrně přehledné a logické, nicméně některé prvky rozhraní mohou být pro uživatele zbytečně matoucí nebo složité. Kupříkladu hned na hlavní stránce kurzu se nachází výrazné tlačítko pro smazání pokroku celého kurzu. Přestože je tato funkce velice důležitá, nejspíše nebude tak často využívána a pozice a výraznost tohoto tlačítka může akorát mást uživatele a odpoutávat jeho pozornost.

Na hlavní stránce kurzu je obsah kurzu rozdělen do seznamu rozbalovacích modulů, v nichž se mohou nacházet různé lekce, kvízy nebo projekty. Struktura kurzu je tak uspořádána přehledně a je snadné se v ní orientovat.

% Odstavec: Navigace v aplikaci
Pro navigaci v aplikaci je využita horní navigační lišta, ve které se nachází rozbalovací možnosti % write - tohle nevím jestli nenazvat jinak
s odkazy na všechny důležité části aplikace. Dále bývá na jednotlivých stránkách využívána levá boční navigační lišta, která slouží k navigaci mezi funkcemi, které jsou relevantní pro danou stránku.

% Odstavec: Srovnání webového a mobilního prostředí

Mobilní aplikace Codecademy Go je výrazně jednodušší, co se funkcí týče. Umožňuje základní funkce zapisování si kurzů a učení se jednotlivých témat. V rámci kurzu jsou k dispozici pouze výrazně zjednodušené lekce, které se liší od lekcí a aktivit ve webové aplikaci. 

% Odstavec: Animace a efekty a mikro interakce
 
V aplikaci jsou animace a efekty použity velice zřídka nebo téměř vůbec. Celá stránka tak působí velice statickým dojmem. 

% Review
\subsubsection{Způsoby učení}

% Struktura kurzu (lineární, tematické moduly, career paths, kapitoly atd.)
Co se struktury kurzů týče, Codecademy nabízí hned několik druhů kurzů. Mezi ně patří kurzy, kariérní cesty a cesty schopností. Všechny tyto druhy kurzů mají stejnou strukturu a liší se zejména svojí obsáhlostí. Dále Codecademy nabízí řadu doplňujících materiálů, kterými jsou různé nástroje pro učení, dokumentace, videa, články a další.

Jednotlivé kurzy jsou rozděleny do jednotlivých modulů. V rámci nich se nachází různé lekce, kvízy, projekty, články nebo další aktivity.

% Jak se spustí učení a jak vypadá study session

Na hlavní stránce kurzu je uživateli zobrazeno tlačítko "Pokračovat v učení", které uživateli spustí lekci nebo aktivitu, u které uživatel naposledy skončil. Jednotlivé lekce typicky skládají z cvičení, která zabírají 5-10 minut. % Write - zde zkontrolovat jak správně psát v češtině 5-10

Cvičení se skládá z textového vysvětlení látky a dále jednotlivých úkolů. Uživatel má k dispozici textový editor, ve kterém může psát a editovat předpřipravený kód. Dále má k dispozici okno, ve kterém je vidět výstup daného kódu.

Kromě zmíněných prvků cvičení si může uživatel také otevřít tzv. tahák se shrnutím daného tématu. Dále si může uživatel spustit přiložené YouTube video, ve kterém učitel prochází daným cvičením a vysvětluje jednotlivé kroky.

Učení je přehledné a uživatel má k dispozici dostatek materiálů pro procvičení a pochopení dané látky. Potenciální nevýhodou může být struktura samotných cvičení. Uživatel je nucen si před každým cvičením přečíst několik stránek textu, což může být pro některé uživatele zdlouhavé a nezáživné.

V rámci učení má uživatel k dispozici AI asistenta, který mu v chatovém okně může vysvětlit danou látku nebo poradit s řešením daného úkolu. 

% Druhy cvičení
V rámci lekcí jsou uživateli zobrazována především cvičení, v nichž uživatel doplňuje a přepisuje předpřipravený kód přímo v editoru nebo píše kód od začátku na základě zadání. 

V rámci kvízů jsou uživateli zobrazovány různé otázky, typicky formou výběru z několika odpovědí nebo výběru zda je daný výrok pravdivý nebo nepravdivý.

Jednou z aktivit jsou takzvané projekty, které svou strukturou připomínají delší cvičení v jednotlivých lekcích. Hlavním rozdílem je, že uživatel nemá možnost nahlédnout do správného řešení a jeho řešení není nijak kontrolováno. Uživatel tak musí daný úkol vyřešit sám bez pomoci. Jednotlivé úkoly pak sám může označit za splněné.  

% Review
\subsubsection{Způsoby motivace uživatele}

% Jaké formy gamifikace aplikace používá
Na rozdíl od aplikace Mimo nepoužívá Codecademy zdaleka tolik forem gamifikace. Aplikace ovšem motivuje uživatele několika následujícími způsoby.

\paragraph{Ukazatele pokroku} Uživatel má jak na domovské stránce, tak na hlavní stránce kurzu k dispozici ukazatele pokroku, které znázorňují, kolik procent daného kurzu uživatel dokončil. 

% Review
\subsubsection{Způsoby monetizace}

% Free tier

% Premium tier


% Free tier


% Review
\subsubsection{Silné a slabé stránky}

% Silné stránky


% Slabé stránky
Mírnou nevýhodou Codecademy může být složitost uživatelského rozhraní. Uživatelé mají k dispozici značné množství materiálů a nástrojů pro učení, nicméně % tady dopsat že to může potenciálně zesložiťovat učení, protože uživatel neví, co používat k učení 

% Review
\subsubsection{Další poznatky}

% Zajímavé funkce

% study plan

% (Zajímavé detaily nebo cokoliv co mě napadne)
